% Generated draft: requires mathematical and source review.
\documentclass[11pt]{article}
\usepackage[margin=25mm]{geometry}
\usepackage{fontspec}
\setmainfont{Latin Modern Roman}
\usepackage{amsmath,amssymb,mathtools,unicode-math}
\setmathfont{Latin Modern Math}
\usepackage{enumitem,needspace,xurl,hyperref}
\setlength{\parindent}{0pt}
\setlength{\parskip}{5pt}
\setlength{\emergencystretch}{3em}
\newcommand{\N}{\mathbb{N}}
\newcommand{\Z}{\mathbb{Z}}
\newcommand{\Q}{\mathbb{Q}}
\newcommand{\R}{\mathbb{R}}
\newcommand{\C}{\mathbb{C}}
\newcommand{\F}{\mathbb{F}}
\newcommand{\A}{\mathbb{A}}
\newcommand{\PP}{\mathbb{P}}
\newcommand{\E}{\mathbb{E}}
\newcommand{\eps}{\varepsilon}
\newcommand{\dd}{\,\mathrm{d}}
\begin{document}
\section*{UnsolvedMath ID: 800194}
\textbf{Source problem code:} catalogue entry
\subsection{Definitions and mathematical statement}
The source asks about the integral quadratic form \(q\) associated with a holomorphic symplectic manifold \(X\). The supplied extract omits the standing hypotheses on \(X\), the normalization of \(q\), and the definition of the Fujiki index \(f_X\) in Section 1.3. Consequently, the following standard interpretation is conditional on confirmation of those omitted conventions.

Under that interpretation, let \(X\) be an irreducible holomorphic symplectic manifold: a compact, simply connected Kähler manifold for which \(H^0(X,\Omega_X^2)\) is generated by an everywhere nondegenerate holomorphic two-form. Write \(\dim_{\C}X=2n\), and let \(L_X=H^2(X,\Z)/\operatorname{Tor}H^2(X,\Z)\). Let \(B_X\) be the Beauville--Bogomolov symmetric bilinear form, with its primitive integral normalization and sign chosen positive on Kähler classes, and write \(q_X(v)=B_X(v,v)\). The lattice \((L_X,B_X)\) is called even if
\[
q_X(v)\in 2\Z\qquad\text{for every }v\in L_X.
\]
Here an isomorphism of lattices means an isomorphism of abelian groups preserving the bilinear forms.

\begin{enumerate}
\item Is \((L_X,B_X)\) even for every \(X\) in the source's intended class?
\item More generally, which isomorphism classes of integral lattices \((L_X,B_X)\) occur as \(X\) varies over that class? The interpretation of ``possibilities for \(q\)'' as classification up to integral lattice isomorphism is suggested by the source's examples, but is not explicitly specified in the extract.
\item Which values can the Fujiki index \(f_X\), with exactly the convention of the source's Section 1.3, take?
\end{enumerate}

For clarity about the last question, the usual Fujiki proportionality coefficient \(c_X\), for a fixed normalization of \(q_X\), is characterized by
\[
\int_X\alpha^{2n}=c_X\,q_X(\alpha)^n
\qquad\text{for all }\alpha\in H^2(X,\C),
\]
where \(q_X\) is extended to complex cohomology. The supplied extract does not specify whether \(f_X=c_X\) or whether a dimension-dependent factor intervenes. Thus this equation does not supply the missing definition of \(f_X\), and the third question cannot be made numerically precise from the extract alone.

The source supplies the following background lattice isomorphisms, attributed there to reference [B1]:
\[
\begin{aligned}
H^2(S^{[r]},\Z)&\cong H^2(S,\Z)\mathbin{\perp}\langle 2-2r\rangle,\\
H^2(K_r(T),\Z)&\cong H^2(T,\Z)\mathbin{\perp}\langle -2-2r\rangle.
\end{aligned}
\]
Here \(\langle m\rangle\) denotes a rank-one lattice generated by a vector of square \(m\), and \(\perp\) denotes orthogonal direct sum. The notation conventionally refers to a Hilbert scheme of points on a K3 surface \(S\) and a generalized Kummer variety associated with a complex torus \(T\); their definitions, the hypotheses on \(T\), and the valid ranges of \(r\) are omitted from the supplied extract and require Section 1.2. These displayed formulas are retained as supplied background, without extending them to unspecified boundary values. The source also states that the lattices of O'Grady's two examples were computed in [R] and are even.

The preceding source paragraph discusses a subalgebra \(\mathcal H\) determined by \(q\) and \(\dim X\), and an \(\mathcal H\)-module \(\mathcal H^\perp\). It states that this module is nonzero for examples (a) and (b) of Section 1.2, except for \(S^{[2]}\) with \(S\) a K3 surface. Those objects and examples are not defined in the extract; this background assertion is preserved as a source assertion and does not add a separate question to the three questions above.
\subsection{Short English statement}
Beauville asks whether the integral Beauville--Bogomolov lattice is always even, which such lattices can occur, and which Fujiki indices can occur. The supplied record labels the problem open. The extract omits the standing hypotheses and normalization conventions, especially the definition of the Fujiki index, so the formulation above remains conditional and requires human mathematical review.
\subsection{Sources}
\begin{itemize}
\item Local UnsolvedMath registry, permanent ID \(800194\), exact source problem code \texttt{catalogue entry}, title ``Holomorphic Symplectic Geometry --- Question 3.'' The supplied record attributes the question to Beauville, with proposed year 2010, and labels its status \texttt{open}; this is a supplied status label, not an independently verified assessment.
\item Beauville, \emph{Holomorphic Symplectic Geometry: A Problem List} (2010), Question 3. Supplied source URL: \url{https://arxiv.org/abs/1002.4321}. The supplied provenance records access on 2026-10-16 and extraction by \texttt{arxiv\_latex}; these are source-record provenance claims, not a claim of access or verification in preparing this draft.
\item The extract attributes the two displayed lattice isomorphisms to [B1] and the computation of the lattices of O'Grady's two examples to [R]. Full bibliographic entries for [B1] and [R] were not supplied. Sections 1.2 and 1.3 are referenced but their contents were not supplied.
\end{itemize}
\end{document}
